%==============================================================================
\section{Correlation Matrix Plot --- Ma trận tương quan}
\label{sec:correlation-matrix}
%==============================================================================

\begin{dinhnghia}[title={Correlation matrix plot}]
Ma trận tương quan là biểu diễn đồ họa của \textbf{tương quan cặp} giữa các biến trong dataset.
Mỗi ô (hoặc heatmap) cho biết hệ số tương quan giữa hai biến; đường chéo thường thể hiện phân phối từng biến (tùy công cụ).
\end{dinhnghia}

\begin{ynghia}[title={Hệ số tương quan}]
Đo mức \textbf{quan hệ tuyến tính} giữa hai biến, trong khoảng $[-1, 1]$:
\begin{itemize}
  \item $+1$: tương quan dương hoàn hảo (biến này tăng thì biến kia tăng).
  \item $-1$: tương quan âm hoàn hảo.
  \item $0$: không có tương quan tuyến tính (trong ý nghĩa hệ số Pearson).
\end{itemize}
\end{ynghia}

\subsection{Ma trận tương quan trong Python}

\begin{phuongphap}[title={Dataset Iris}]
Ví dụ dùng Iris (scikit-learn): 150 mẫu, 4 feature, 3 loài.
\end{phuongphap}

\subsection{Ví dụ --- Heatmap tương quan}

\begin{vidu}[title={Code đầy đủ (Pandas + Seaborn)}]
\begin{lstlisting}
import numpy as np
import pandas as pd
import seaborn as sns
import matplotlib.pyplot as plt
from sklearn.datasets import load_iris

iris = load_iris()
data = pd.DataFrame(iris.data, columns=iris.feature_names)
target = iris.target

corr = data.corr()
sns.set(style='white')
mask = np.zeros_like(corr, dtype=bool)
mask[np.triu_indices_from(mask)] = True
f, ax = plt.subplots(figsize=(11, 9))
cmap = sns.diverging_palette(220, 10, as_cmap=True)
sns.heatmap(corr, mask=mask, cmap=cmap, vmax=.3, center=0,
   square=True, linewidths=.5, cbar_kws={"shrink": .5})
plt.show()
\end{lstlisting}

\texttt{corr()} trên DataFrame; \texttt{mask} che tam giác trên để tránh lặp; \texttt{heatmap} tô màu theo hệ số.
\end{vidu}

\begin{ynghia}[title={Đọc biểu đồ}]
\begin{itemize}
  \item \texttt{sepal width} và \texttt{petal length}: tương quan âm vừa ($\approx -0{,}37$).
  \item \texttt{petal length} và \texttt{petal width}: tương quan dương mạnh ($\approx 0{,}96$).
  \item \texttt{sepal length} và \texttt{petal length}: tương quan dương ($\approx 0{,}87$).
\end{itemize}
\end{ynghia}

\begin{vidu}[title={Output}]
\tsimg{05_correlation_matrix_plot/img01.jpg}
\end{vidu}
